\documentclass[12pt]{article}

\usepackage{caption}
\usepackage{subcaption}
\usepackage{authblk}

\usepackage{amsmath}
\usepackage{amsfonts}
\usepackage{url}
\usepackage{amsthm}
\usepackage{color}
%\usepackage{stmaryrd}
\usepackage{graphicx}
\usepackage{float}
\newtheorem{mydef}{Definition}
\newtheorem{mylem}{Lemma}
\newtheorem{myass}{Assumption}
\newtheorem{myprop}{Proposition}
\newtheorem{mycor}{Corollary}
\newtheorem{myrem}{Remark}
\newtheorem{mythm}{Theorem}

\newcommand{\bs}[1]{{\boldsymbol #1}}
\newcommand{\hot}[1]{\textcolor{red}{#1}} 
\newcommand{\bsxi}{\bs \xi}

\usepackage[margin=1in]{geometry}




\title{Steady-State Analysis of Gravitational Effects on Hemodynamics}
\author[1,2]{Zan Ahmad}
\author[]{Karen Ong}
\author[]{Charles Peskin}
\author[]{Rebecca Blue}
\author[]{Alanna Kennard}
\affil[1]{Courant Institute of Mathematical Sciences, New York University}
\affil[2]{Johns Hopkins University Department of Applied Mathematics and Statistics}
\date{December 2022}

\begin{document}

\maketitle
\begin{abstract}
    In this paper, we derive solutions for a linear model of the circulation with the incorporation of gravity. Key features of this approach include feedback control elements to model the autonomic nervous system's regulation of heart rate and reserve volume, compartmentalization of the upper and lower circulation to allow for gravitational acceleration in the vertical direction to be incorporated, and the inclusion of partial collapse of the systemic veins. We utilize a time-independent steady state modeling approach, the simplicity of which enables us to predict an individual's tolerance to high gravitational acceleration in a computationally inexpensive manner. The parameters of the model were calibrated with a subject's personalized quantitative centrifuge simulation data and physical measurements and results of the model were compared to qualitative information about the subject's tolerance to exposure to hypergravity. 
\end{abstract}

\section{Introduction}
The biophysical processes underlying hemodynamics and the regulation of blood flow are complex and nonlinear. They involve the pulsatile contractions of the heart, which are regulated by the pacemaker in the sino-atrial node in healthy individuals. The rate of these contractions is in turn mediated by baroreceptors located in the carotid sinus and aortic arch. The pulsatile contractions of the heart, smooth muscle and skeletal muscle, along with the compliance of blood vessels, regulate the flow of fluid throughout the vascular system. Several phenomena are known to cause temporary changes in blood flow, causing reduced blood flow to the brain and resulting in a loss or alteration of consciousness. These include orthostatic hypotension,  transient ischemic attack, loss of blood volume\textcolor{red}{(cite)}. 
\par Of particular interest to us is gravity-induced loss of consciousness (G-LOC), which occurs when the sudden increase in the vertical acceleration, $G_{\text{z}}$, causes blood to pool in the lower extremities. To investigate the relationship between high $G_\text{z}$ and hyper-gravity induced pathologies such as G-LOC, we implement a circulation model which incorporates the effects of gravitational acceleration. While the effect that vertical and lateral acceleration has on the body is fundamentally non-linear, we postulate that for the binary prediction task of an individual's G-tolerance, it can be approximated efficiently by a linear model for the control of blood flow. This model has the advantage of being computationally inexpensive and readily interpretable.
\par Many works have already proposed mathematical descriptions of the circulatory system. Models typically follow the analog circuit model put forth by Guyton. Many models are ordinary differential equation (ODE) or partial-differential equation (PDE) based. These models typically simulate the heart using an pulsatile integral model \textcolor{blue}{elaborate} and the flow throughout the systemic system using using an inflow/outflow model for each subsystem. The inflow/outflow to a particular organ can be minimalistically modeled \cite{diaz2019computational} following Guyton's theoretical framework. Arteries and veins can also be modeled more precisely using deformable cyclinder models and flow through them using a valve model at each artery or vein junction.\cite{zhang2017gravitational}.

These models are often able to simulate accurate waveforms for \textcolor{red}{check} pressure and flow. They are then run to a steady-state to measure \textcolor{red}{complete.}. However many of these multiscale model formulations are not parametrizable in practice (either the number of parameters is overly large or the measures required too invasive) or amenable to predicting practical quantities. For instance Guyton's original model involved intra-atrial and intra-arterial manometers to measure blood pressure (invasive), Zhang's model involved measuring all \textcolor{red}{50 arteries and 120 veins in the human body} (impractical) and the Mohammadyari and Artiles models use many constants that have not been selected by rigourous statistical analysis.
\textcolor{red}{to complete}
Zhang et al \cite{zhang2017gravitational}, Artiles-Diaz \cite{diaz2019computational} and  Mohammadyari \cite{mohammadyari2021modelling} follow Guyton et al, by using an analog circuit model of the human circulation. However they do so using three different approaches.


Studies on the effects of centrifugation-induced hypergravity typically focus on post-facto gravity-induced changes on the circulatory system just succeding centrifufation. While we are unsure of the time duration required for the steady state evolved during centrifugation to revert to the steady state under normal G, these studies provide quantitative confirmation of the hypothesized effects of Gx on blood pressure. Serrador et al spin subjects at 3Gz for 30 minutes and then measure their eye-blood pressure and cerebral perfusion pressure. Both decreased as a result of Gz exposure (no anti-G straining maneuvers were employed).
Ossard et al \cite{ossard1996cerebral} measure that mean blood flow velocity significantly decreases at 3-4Gz (by ---) before subjects performed the valsalva maneuver.

\section{Model Derivation}
The foregoing formulations describe a time-independent, steady-state flow model which incorporates the effects of gravity, $G_\text{z}$ on haemodynamics. Our model descriptions are comprised of two parts: (i) a derivation of the steady-state blood flow model and (ii) incorporation of an idealized controller of heart rate.  

\subsection{Blood Flow Model}
The blood flow model used here describes the circulation as a series of compartments that are either compliance chambers or resistance elements. Each compartment is identified with a unique index $i$. We assume the following relation between compliance $C_i$, pressure $P_i$, and volume $V_i$:
\begin{align}
V_i = V^0_{i} + C_i P_i = V^0_{i} + C_i (P^{ \text{interior}}_i - P^{ \text{exterior}}_i).
\label{eq:comp_chamber}
\end{align}

Compliance is the ability of a structure to change shape as a function of volume and pressure. The contractions of cardiac chambers results in the compliance periodically oscillating between a minimum and maximum value during the phases of systole, when the chamber contracts, and diastole, when the chamber relaxes, respectively. The dead volume $V^0_{i}$ is equal to the residual chamber volume at zero pressure. We will interchangeably refer to dead volume as reserve volume, since a reduction in $V^0$ may compensate for a loss in total volume, which we define as the sum of the volumes in each compartment: 
\begin{align}
    V_\text{total} = \sum_i V_i \label{eq:Vtotal}
\end{align}
and similarly, the total reserve volume is given by the sum of the dead volumes in each compartment. 
\begin{align}
    V^0_\text{total} = \sum_i V^0_i \label{eq:V0total}
\end{align}
Subtracting equation (\ref{eq:V0total}) from (\ref{eq:Vtotal}) and substituting (\ref{eq:comp_chamber}) yields:
\begin{align}
    V_\text{total} - V^0_\text{total} = \sum_iC_iP_i \label{eq:V0-V}
\end{align}

\begin{figure}[h!]
    \centering
    \includegraphics[scale=0.3]{model-schematic.png}
    \caption{Schematic of the compartmentalization of the circulation used in this model. (u = upper, l = lower, p = pulmonary, s = systemic, pv = pulmonary veins, sv = systemic veins, sa = systemic arteries, pa = pulmonary arteries, RA = right atrium, RV = right ventricle, LA = left atrium, LV = left ventricle. }
    \label{fig:circ-diagram}
\end{figure}

In a resistance vessel, the flow is proportional to the pressure difference from one end of the vessel to the other. Thus, for any flow through
a resistance vessel between chamber $i$ and $j$, $Q_{ij}$ is given by Ohm's Law/Poiseuille's Equation: 
\begin{align}
    Q_{ij} = \dfrac{P_{i} - P_j}{R_{ij}} \label{eq:Ohm}
\end{align}

In the foregoing formulations, we adopt a continuous flow view point, which means that we are averaging over the cardiac cycle. Moreover, we consider a steady-state situation, where the pressures, flows, and volumes are all time-independent. We model the systemic arteries (sa), systemic veins (sv), pulmonary arteries (pa) and pulmonary veins (pv) as compliance vessels and tissues and organs as resistance vessels. Furthermore, the systemic circulation is further compartmentalized into upper and lower components. There is also a thoracic compartment to our model which will be described in the discussion on partial venous collapse. See figure \ref{fig:circ-diagram} for schematic of the connectivity of the described elements.

Cardiac output or systemic flow is defined as the volume of blood delivered to the systemic circulation over time and can be computed by: 
\begin{align}
    Q_\text{sa,sv} = FV_{\text{stroke}} \label{eq: CO}
\end{align}
where $Q_\text{sa,sv}$ is the flow between the systemic arteries and systemic veins, $F$ is heart rate in beats per unit time, and $V_{\text{stroke}}$ is the volume of blood ejected per heart beat. In the systemic circulation, there are multiple tissues and organs connected in parallel; their flows add to give the total systemic flow:
\begin{align}
    Q_\text{sa,sv} = \sum_j Q_{\text{s}_j} \label{eq:sum}
\end{align}

We make the steady state assumption that flow entering a vessel is equal to flow exiting a vessel, and thus, we may drop the subscript on $Q_{ij}$ for any $i$ and $j$ and simply refer to $Q$ as flow in general. 

Since $F$ is the same for both the left and right side of the heart (as a result of there being a single pacemaker, known as the sino-atrial node) the stroke volumes of the left and right heart are equal and are approximated by using equation (\ref{eq:comp_chamber}) for the volume in the left ventricle at the end of diastole: 
\begin{align}
    V_\text{stroke} = C_\text{LVD}(P_\text{RA}-P_\text{thorax})\label{eq: Vstroke-left}\\= C_\text{RVD}(P_\text{pv}-P_\text{thorax}) \label{eq: Vstroke-right} 
\end{align}
where $C_\text{LVD}$ and $C_\text{RVD}$ are the left and right ventricular diastolic compliances respectively. Combining equations (\ref{eq: CO}) with (\ref{eq: Vstroke-left}) and (\ref{eq: Vstroke-right}), respectively, yields the following: 
\begin{align}
    Q = FC_\text{LVD}(P_\text{RA}-P_\text{thorax}) \label{eq: Q}\\ 
      = FC_\text{RVD}(P_\text{pv}-P_\text{thorax}) \label{eq: heart rate}
\end{align}

Additionally, by equation (\ref{eq:Ohm}) and the steady state assumption, we have pulmonary flow given by:
\begin{align}
    Q = \frac{P_\text{pa}-P_\text{pv}}{R_\text{p}} \label{eq:pulmflow}
\end{align}
\subsection{Incorporation of Gravity}
To incorporate the effects of gravity in the vertical direction, we compartmentalize the systemic circulation into an upper and lower component with their own pressure, flow and resistance values. Then, by equation (\ref{eq:Ohm}), 
\begin{align}
    Q^\text{u}_\text{s} = \frac{P^\text{u}_\text{sa} - P^{\text{u}}_\text{sv}}{R^\text{u}_\text{s}} \label{eq:sysflowupper}\\
    Q^\text{l}_\text{s} = \frac{P^\text{l}_\text{sa} - P^{\text{l}}_\text{sv}}{R^\text{l}_\text{s}} \label{eq:sysflowlower}
\end{align}
Note that by equation (\ref{eq:sum}), 
\begin{align}
    Q = Q^\text{u}_\text{s} + Q^\text{l}_\text{s} 
\end{align}

Then, we are able consider the blood between the upper and lower compartment of the systemic circulation as a fluid-filled column and use Bernouli's equation to obtain the following relationship: 
\begin{align}
    P^\text{l}_{\text{sa}} - P^\text{u}_{\text{sa}} = \rho g(H^\text{u} - H^\text{l}) \label{eq:Pressure_diff_height}
\end{align}
where, $\rho$ is the density of blood, $g$ is the gravitational acceleration, and $H^\text{u}$ and $H^\text{l}$ denote upper and lower heights respectively such that 
\begin{align}
    H^\text{u}> 0 > H^\text{l}.
\end{align}
\subsection{Partial Venous Collapse}
For most of the systemic circulation, the external pressure is that of the atmosphere, which we denote as zero pressure. For the pulmonary circulation and both sides of the heart, the external pressure is the intrathoracic pressure, denoted as $P_\text{thorax}$. $P_\text{thorax}$ varies during the cycle of breathing and is lower during inspiration and higher during expiration. On average, the value of $P_\text{thorax}$ is negative and is to be thought of as the pressure in the tissue surrounding the lungs. The lungs are elastic and tend to collapse and a pressure difference between the air in the lungs and the air in the surrounding tissue is needed to keep them open. Negative intrathoracic pressure is generated by contraction of the diaphragm. The heart chambers are lumped into one chamber which we choose to be the right atrium. Depending on right atrial pressure, $P_\text{RA}$, we get the following equations
\begin{align}
    P_{\text{sv}}^\text{u} &= \text{max}\{0,P_{\text{RA}}-\rho g H^\text{u}\} \label{eq:cond_PsvU} \\
    P_{\text{sv}}^\text{l} &= \text{max}\{(0,P_{\text{RA}})-\rho g H^\text{l}\}\label{eq:cond_PsvL}
\end{align}
which give rise to three cases. 

Case I:
\begin{align}
    P_\text{thorax}<P_\text{RA}<0
\end{align}



Case II: 
\begin{align}
    0 < P_\text{RA} < \rho g (-H^\text{l}).
\end{align}
In this case, $P^\text{u}_\text{sv}=0$ and $P^\text{l}_\text{sv}= P_\text{RA} + \rho g (-H^\text{l})$.

In cases I and II, we get a localized partial collapse of the systemic veins, just upstream of their entry into the thoracic compartment. This is referred to as "partial venous collapse". 

Lastly, Case III: 
\begin{align}
   \rho g (H^\text{u})<   P_\text{RA}.
\end{align}
In this case, $P^\text{u}_\text{sv}= P_\text{RA} - \rho g (-H^\text{u})$ and $P^\text{l}_\text{sv}= P_\text{RA} + \rho g (-H^\text{l})$ and here we have no collapse of the systemic veins. 

We now have a linear system of equations that we can solve analytically for each of the three cases. In the uncontrolled case, our parameters are the total reserve volume, compliance and resistance values of the arteries and veins, stroke volume, the upper and lower compartment heights, total volume, gravitational acceleration, blood density, and heart rate. The need for feedback control can be clearly seen by considering the example of what occurs during exercise. When we begin to exercise, the blood vessels in our systemic circulation dilate to increase blood flow to the tissues that are being depleted of oxygen. This dilation results in the radius of the vessel increasing and thereby reducing the systemic resistance. Clearly, by equations (\ref{eq:sysflowupper}) and (\ref{eq:sysflowlower}), this will cause the systemic arterial pressure to drop. Because heart rate is a fixed parameter, the cardiac output, $Q$, will not increase sufficiently. Baroreceptors detect this drop in blood pressure from a set value and adjust the heart rate accordingly. 

\subsection{Idealized Controller}

In this subsection, we attempt to remedy the issues of the uncontrolled circulation by considering $F$ and $V^0_\text{total}$ to be unknown variables rather than parameters, so that $Q$ can be proportionally adjusted during dynamic changes to the circulation. Since total volume is still fixed, the reduction of total reserve volume serves to increase the volume of distributed blood between the chambers, thereby increasing flow. 
We assume that the autonomic nervous system is able to make adjustments to heart rate and total reserve volume to maintain specific target values of $P_\text{sa}^\text{u}$ and also the pressure difference that stretches the wall of the right atrium, $\Delta P_\text{RA}$. Thus, we have new parameters $\left( P_{\text{sa}}^\text{u} \right)^*$ and $\left( \Delta P_{\text{RA}} \right)^*$ defined as
\begin{align}
    P_{\text{sa}}^\text{u} &= \left( P_{\text{sa}}^\text{u} \right)^* \label{eq:crtlled_Psa}\\
    P_{\text{RA}} - P_{\text{thorax}} &= \left( \Delta P_{\text{RA}} \right)^* \label{eq:crtlled_Pra_pth}
\end{align}

For the model to be valid, $V_\text{total}^0$ must be nonnegative as a working system always has some venous reserve volume. We will derive analytical solutions for $F$ and $V^0_\text{total}$ for each case separately, with the following procedure: 
\begin{enumerate}
    \item Identify the case defining condition for validity on $P_\text{thorax}$.
    \item Use the upper systemic venous pressures and the control equations (\ref{eq:crtlled_Psa}) and (\ref{eq:crtlled_Pra_pth}) to solve for the lower systemic venous pressure.  Solve for the upper and lower systemic flows with the appropriate computed pressure values to get $Q$ and $F$ as a function of parameters. 
    \item Solve for pulmonary pressures to express total reserve volume $V^0_\text{total}$ as a function of parameters.
\end{enumerate}

\subsubsection{Case I}
Since \(P_{\text{RA}}=P_{\text {thorax }}+\left(\Delta P_{\text{RA}}\right)^{*}\), the condition for validity of case I can  be stated as follows 

\begin{align}
    P_{\text {thorax }} \leq-\left(\Delta P_{\text{RA}}\right)^{*}
\end{align}
\vspace*{-\baselineskip}

In this case, we define the upper and lower systemic venous pressures as 
\begin{align}
    P^\text{u}_\text{sv}&=0  \\
    P^\text{l}_\text{sv}&=\rho g (-H^\text{l})
\end{align}
\( P_\text{sv}^\text{u} =0 \) and \( P_\text{sv}^\text{l} =\rho g-H^\text{l}) \).
Substituting the controlled value for $P_{\text{sa}}^\text{u}$ into equation (\ref{eq:Pressure_diff_height}) yields 
\begin{align}
    P_\text{sa}^\text{l} &=\left(P_{\text{sa}}^\text{u}\right)^{*}+\rho g\left(H^\text{u}+\left(-H^{\text{l}}\right)\right) \label{eq:PsaL_caseI_sol}
\end{align}
Using equations (\ref{eq:sysflowupper}), (\ref{eq:sysflowlower}) and (\ref{eq:PsaL_caseI_sol}), the systemic pressures specify the flows 
\begin{align}
Q_\text{s}^\text{u} &=\frac{1}{R_\text{s}^\text{u}}\left(P_{\text{sa}}^\text{u}\right)^{*} 
\end{align}
% \vspace{-\baselineskip}
\begin{align}
\begin{split}
    Q_\text{s}^\text{l} &= {} \frac{1}{R_\text{s}^\text{l}}\left(P_\text{sa}^\text{l}-P_\text{sv}^\text{l}\right) \\
& = \frac{1}{R_\text{s}^\text{l}}\left(\left(P_{\text{sa}}^\text{u}\right)^{*}+\rho g H^\text{u}\right)
\end{split}
\end{align}
With both upper and lower systemic flows known, the cardiac output is simply their sum: 
\begin{align}
  Q &=  \left(\frac{1}{R_\text{s}^\text{u}}+\frac{1}{R_\text{s}^\text{l}}\right)\left(P_{\text{sa}}^\text{u}\right)^{*}+\frac{\rho g H^\text{u}}{R_\text{s}^\text{l}}     \label{eq:Q-case1}
\end{align}
Then, by using equation (\ref{eq: heart rate}) and substituting (\ref{eq:crtlled_Pra_pth}) and (\ref{eq:Q-case1}), an equation for heart rate as a function of parameters can be obtained: 
\begin{align}
 F &=\frac{Q}{C_\text{RVD}\left(P_{\text{RA}}-P_{\text {thorax }}\right)} \label{eq: heartrate2} \\
&=\frac{\left(\frac{1}{R_\text{s}^\text{u}}+\frac{1}{R_\text{s}^\text{l}}\right)\left(P_{\text{sa}}^\text{u}\right)^{*}+\frac{\rho g H^\text{u}}{R_\text{s}^\text{l}}}{C_\text{RVD}\left(\Delta P_{\text{RA}}\right)^{*}}
\end{align}



% not adding eq 44-47 as pulmonary pressure not an output of interest to the community \textcolor{red}{ add 44-47 ?} I included the pumonary pressures because they are part of the derivation that leads to the output of the heart rate and reserve volume. -Zan 
The pulmonary pressures can be solved for by setting equations (\ref{eq: Q}) and (\ref{eq: heart rate}) equal to each other and substituting equation (\ref{eq:crtlled_Pra_pth}):
\begin{align}
    C_\text{LVD}(P_\text{pv} - P_\text{thorax}) &= C_\text{RVD} \left( \Delta P_\text{RA} \right)^* \\
    P_\text{pv} - P_\text{thorax} &= \frac{C_\text{RVD}}{C_\text{LVD}}\left( \Delta P_\text{RA} \right)^* 
\end{align}
We know from equation (\ref{eq:pulmflow}) that $P_\text{pa}-P_\text{pv} = QR_\text{p}$ and adding this quantity to both sides results in: 
\begin{align}
    P_\text{pa} - P_\text{thorax}= \frac{C_\text{RVD}}{C_\text{LVD}}\left( \Delta P_\text{RA} \right)^* + QR_\text{p}
    \end{align}
Using equation (\ref{eq:V0-V}) and our case I assumptions, we also arrive at an expression for total reserve volume as a function of parameters:
% multi line version of vto solution
% \begin{align}
% \begin{split}
%   V_{\text {total }}^{0} &= {} V_{\text {total }} -C_{p} \frac{C_{R V D}}{C_{L V D}} \left(\Delta P_{\text{RA}} \right)^{*} -\left(T_{p} G_{s} +C_{sa}\right) \left(P_{\text{sa}}^\text{u}\right)^{*}  \\
%   & -\left(T_{p} G_{s}^{l} +C_\text{sa}^{l}\right) \rho g H^{u} \\
%  & -C_{s}^{l} \rho g\left(-H^{l}\right)  
% \end{split}
% \end{align}

\begin{align}
V_{\text{total}}^0 &= V_\text{total}- C_\text{p} \frac{C_\text{RVD}}{C_\text{LVD}}\left (\Delta P_{\text{RA}}\right)^* -\left(T_\text{p} G_\text{s}+C_\text{sa}\right)\left(P_{sa}^\text{u}\right)^* - \left(\dfrac{T_\text{p} }{R_\text{s}^\text{l}} +C_\text{sa}^L\right) \rho g H^\text{u} -C_\text{s}^\text{l} \rho g-H^\text{l}) \label{eq:solution_VTO_caseI}
\end{align}
where the arterial compliance is the sum of the compartment compliances: 
\begin{align}
    C_\text{sa} =C_\text{sa}^{\text{u}}+C_\text{sa}^{\text{l}},
\end{align} 
the conductance term $G_{\text{s}}$ is the sum of the upper and lower conductances (reciprocal resistances):
\begin{align}
G_\text{s} = G_\text{s}^\text{u} + G_\text{s}^\text{l} = \frac{1}{R_\text{s}^\text{u}}+\frac{1}{R_\text{s}^\text{l}},
\end{align}
the pulmonary time constant $T_\text{p}$ is 
\begin{align}
    T_\text{p} = C_\text{pa}R_\text{p}
\end{align}
 and the pulmonary compliance, $C_\text{p}$ is
 \begin{align}
     C_\text{p}=C_\text{pa}+C_\text{pv}.
 \end{align}
 
% (and individual \textcolor{red}{resistance} terms are of the form: \(G_{s}^{l}=\frac{1}{R_{s}^{l}} \)
% labeled gs and cs equations
% where \begin{align}
% G_s &= \frac{1}{R_s^u}+\frac{1}{R_s^l}
% \end{align} and following compartment homogenity, \begin{align}
%     C_\text{sa} &=C_\text{sa}^u+C_\text{sa}^l
% \end{align}

\subsubsection{Case II}
In a similar manner, case II which is defined by the following inequality: 
\begin{align}
    0 < P_\text{RA} < \rho g H^\text{u} 
\end{align}
which can be rewritten by equation (\ref{eq:crtlled_Pra_pth}) as
\begin{align}
    -\left( \Delta P_{\text{RA}} \right)^*<P_\text{thorax}<\rho g H^\text{u}.
\end{align}

Since $P_\text{RA}<\rho g H^\text{u}$, there is partial collapse of the upper systemic veins (not extending all the way down to the heart). Thus, similar to case I, we have that $P_\text{sv}^\text{u} = 0$, but because right atrial pressure is nonnegative, there is no systemic venous collapse below the heart, yielding the following: 
\begin{align}
    P_\text{sv}^\text{l} &= \rho g (-H^\text{l}) + P_\text{RA}\\
    &= \rho g (-H^\text{l}) + P_\text{thorax} +\left( \Delta P_{\text{RA}}\right)^*
\end{align}

The upper and lower systemic arterial pressures are the same as in case I, given by equations (\ref{eq:crtlled_Psa}) and (\ref{eq:PsaL_caseI_sol}), and the upper and lower systemic flows follow: 
\begin{align}
    Q_\text{s}^\text{u} &= \frac{1}{R_\text{s}^\text{u}}\left((P_\text{sa}^\text{u})^*+\rho g (H^\text{u} + -H^\text{l}) \right) \\
    Q_\text{s}^\text{l} &= \frac{1}{R_\text{s}^\text{l}}\left((P_{\text{sa}}^\text{u})^{*}+\rho g\left(H^\text{u}+(-H^{\text{l}})\right) - \rho g (-H^\text{l}) - P_\text{thorax} - \left( \Delta P_\text{RA} \right)^*\right)\\
    &= \frac{1}{R_\text{s}^\text{l}}\left( (P_\text{sa}^\text{u})^*+\rho g H^\text{u} - P_\text{thorax} - \left( \Delta P_\text{RA} \right)^*\right)
 \end{align}

Adding the upper and lower systemic flows together yields the cardiac output: 
\begin{align}
Q = \left(\frac{1}{R_\text{s}^\text{u}} + \frac{1}{R_\text{s}^\text{u}}\right) (P_\text{sa}^\text{u})^* + \frac{1}{R_\text{s}^\text{l}} \left(\rho g H^\text{u} - P_\text{thorax} - \left( \Delta P_\text{RA} \right)^*\right)
\end{align}

which in turn allows us to solve for heart rate by equation (\ref{eq: heart rate}) and (\ref{eq:crtlled_Pra_pth}):
\begin{align}
    F &= \frac{Q}{C_\text{RVD}} \left( \Delta P_\text{RA} \right)^* \label{eq:heartrate3}\\
    &= \frac{\left(\frac{1}{R_\text{s}^\text{u}} + \frac{1}{R_\text{s}^\text{u}}\right) (P_\text{sa}^\text{u})^* + \frac{1}{R_\text{s}^\text{l}} \left(\rho g H^\text{u} - P_\text{thorax} - \left( \Delta P_\text{RA} \right)^*\right)}{C_\text{RVD}} \left( \Delta P_\text{RA} \right)^* \label{eq: F-case2}
\end{align}
The pulmonary pressures can be solved by setting equations (\ref{eq: Q}) and (\ref{eq:heartrate3}) equal to each other after solving for :
\begin{align}
    C_\text{LVD}(P_\text{pv} - P_\text{thorax}) &= C_\text{RVD} \left( \Delta P_\text{RA} \right)^* \\
    P_\text{pv} - P_\text{thorax} &= \frac{C_\text{RVD}}{C_\text{LVD}}\left( \Delta P_\text{RA} \right)^* \label{eq:Ppv-Pthorax}
\end{align}
We know from equation (\ref{eq:pulmflow}) that $P_\text{pa}-P_\text{pv} = QR_\text{p}$ and adding this quantity to both sides results in: 
\begin{align}
    P_\text{pa} - P_\text{thorax}= \frac{C_\text{RVD}}{C_\text{LVD}}\left( \Delta P_\text{RA} \right)^* + QR_\text{p} \label{eq:Ppa-Pthorax}
    \end{align}
\begin{align}
    = \frac{C_\text{RVD}}{C_\text{LVD}}+R_\text{p}\left(\left(\frac{1}{R_\text{s}^\text{u}} + \frac{1}{R_\text{s}^\text{u}}\right) (P_\text{sa}^\text{u})^* + \frac{1}{R_\text{s}^\text{l}} \left(\rho g H^\text{u} - P_\text{thorax} - \left( \Delta P_\text{RA} \right)^*\right)\right)
\end{align}
Following the same procedure from case I, we can solve for $V^0_\text{total}$ in terms of parameters using equation (\ref{eq:V0-V}) and our case II assumptions:
\begin{equation}
    \begin{aligned}
    V_\text{total}^0 =  V_\text{total}- C_\text{p} \frac{C_\text{RVD}}{C_\text{LVD}}\left (\Delta P_{\text{RA}}\right)^* -\left(T_\text{p} G_\text{s}+C_\text{sa}\right)\left(P_\text{sa}^\text{u}\right)^* -(T_\text{p}G_\text{s}^\text{l}\\ 
     + C_\text{sa}^\text{l})\rho g H^\text{u} - C_\text{s}^\text{l} \rho g (-H^\text{l})-(C_\text{sv}^\text{l} - T_\text{p}G_\text{s})(P_\text{thorax} + (\Delta P_\text{RA})^*)\label{V0-case2}
    \end{aligned} 
\end{equation}
\subsubsection{Case III}
Repeating the procedure from the previous two cases, we can derive analytical solutions for heart rate and total reserve volume in case III where there is no venous collapse. The defining inequality for this case is 
\begin{align}
    \rho g (H^\text{u})<   P_\text{RA} = P_\text{thorax} + \left( \Delta P_\text{RA} \right)^*
\end{align}
which subsequently gives the condition for validity on $P_\text{thorax}$:
\begin{align}
    \rho g H^\text{u} - \left( \Delta P_\text{RA} \right)^* \leq P_\text{thorax}
\end{align}
The upper and lower systemic venous pressures are defined as:
\begin{align}
   P^\text{u}_\text{sv}&= P_\text{RA} - \rho g (-H^\text{u}) = P_{\text{thorax}} + \left(\Delta P_\text{RA} \right)^* - \rho g (-H^\text{u})\\
   P^\text{l}_\text{sv}&= P_\text{RA} + \rho g (-H^\text{l}) = P_{\text{thorax}}+ \left(\Delta P_\text{RA} \right)^* + \rho g (-H^\text{l})
\end{align}
The upper and lower systemic arterial pressures are the same as the other two cases as given by given by equations (\ref{eq:crtlled_Psa}) and (\ref{eq:PsaL_caseI_sol}) and combining them with equations  (\ref{eq:sysflowupper}) and (\ref{eq:sysflowlower}) result in: 
\begin{align}
\begin{split}
    Q_\text{s}^\text{u} &= \frac{1}{R_\text{s}^\text{u}}\left((P_\text{sa}^\text{u})^*-(P_\text{thorax} +(\Delta P_\text{RA})^* -\rho g H^\text{u}) \right) \\
    &= \frac{1}{R_\text{s}^\text{l}}\left((P_\text{sa}^\text{u})^* + \rho g H^\text{u} - (P_\text{thorax}) - (\Delta P_\text{RA})^*\right) 
\end{split}
\end{align}

\begin{align}
    \begin{split}
    Q_\text{s}^\text{l} = \frac{1}{R_\text{s}^\text{l}}\left((P_\text{sa}^\text{u})^* + \rho g(H^\text{u} + -H^\text{l}))-\\
    (P_\text{thorax} +(\Delta P_\text{RA})^* + \rho g (-H^\text{l})\right))
    \\
    = \frac{1}{R_\text{s}^\text{l}}\left((P_\text{sa}^\text{u})^* + \rho g H^\text{u} - (P_\text{thorax}) - (\Delta P_\text{RA})^*\right) 
    \end{split}
\end{align}
Since the upper and lower pressure differences are the same, the sum of the upper and lower systemic flows results in the following equation for cardiac output: 
\begin{align}
    Q = \left( \frac{1}{R_\text{s}^\text{u}} + \frac{1}{R_\text{s}^\text{l}}\right)\left((P_\text{sa}^\text{u})^* + \rho g H^\text{u} - (P_\text{thorax}) - (\Delta P_\text{RA})^*\right) \label{eq: Q case 3}
\end{align}
Then, heart rate is given in the same way by  rearranging equation (\ref{eq: heart rate}) and plugging in $Q$ from equation (\ref{eq: Q case 3}):
\begin{align}
    F = \frac{\left( \frac{1}{R_\text{s}^\text{u}} + \frac{1}{R_\text{s}^\text{l}}\right)\left((P_\text{sa}^\text{u})^* + \rho g H^\text{u} - (P_\text{thorax}) - (\Delta P_\text{RA})^*\right)}{C_\text{RVD}\left(\Delta P_\text{RA} \right)^*} \label{eq:F-case3}
\end{align}
Then, we use the pulmonary pressures from (\ref{eq:Ppv-Pthorax}) and (\ref{eq:Ppa-Pthorax}) and plug in our case III value for $Q$ from equation (\ref{eq: Q case 3}) into the latter: 
\begin{align}
    P_\text{pa} - P_\text{thorax}= \frac{C_\text{RVD}}{C_\text{LVD}}\left( \Delta P_\text{RA} \right)^* + \left( \frac{1}{R_\text{s}^\text{u}} + \frac{1}{R_\text{s}^\text{l}}\right)\left((P_\text{sa}^\text{u})^* + \rho g H^\text{u} - (P_\text{thorax}) - (\Delta P_\text{RA})^*\right) R_\text{p} 
    \end{align}
Finally, we can write the total reserve volume for case III with the following familiar expression from equation (\ref{eq:V0-V}): 
    \begin{equation}
    \begin{aligned}
    V_\text{total}^0 =  V_\text{total}- C_\text{p} \frac{C_\text{RVD}}{C_\text{LVD}}\left (\Delta P_{\text{RA}}\right)^* -\left(T_\text{p} G_\text{s}+C_\text{sa}\right)\left(P_\text{sa}^\text{u}\right)^* -(T_\text{p}G_\text{s}^\text{l} +\\ 
      C_\text{sa}^\text{l} - C_\text{sv}^\text{u})\rho g H^\text{u} - C_\text{s}^\text{l} \rho g (-H^\text{l})-(C_\text{sv}^\text{l} - T_\text{p}G_\text{s})(P_\text{thorax} + (\Delta P_\text{RA})^*)\label{eq:V0-case3}
    \end{aligned} 
\end{equation}
The borderline between case II and III is defined by
\begin{align}
    P_\text{thorax} + (\Delta P_\text{RA})^*= \rho g H^\text{u}.
\end{align}
Setting equations (\ref{eq: F-case2}) and (\ref{eq:F-case3}) equal gives:
\begin{align}
    \frac{G_\text{s}(P_\text{sa}^\text{u})^*}{C_\text{RVD}(\Delta P_\text{RA})^*}
\end{align}
for both cases. 
Setting the total reserve volume equations for these two cases, (\ref{V0-case2}) and (\ref{eq:V0-case3}), equal and noticing that they only differ with the terms $\rho g H^\text{u}$ and $P_\text{thorax} + (\Delta P_\text{RA})^*$. Thus, heart rate and total reserve volume are continuous functions in the transitions between case II and III. 
\end{document}